\begin{afterword}
\summaryhead

The essay argues that ``everything is uncertain'' or ``everyone is imperfect'' does not make all beliefs or all people equal. Nothing is pure white or pure black, but some grays are much lighter than others, and the difference matters.

The author calls the opposite view the Fallacy of Gray. He illustrates it with a man who shrugged off enormous odds against his view, with bloggers who said that no one can eliminate bias, and with an economist who said that theory is never applied without personal judgment. His reply to each is the same: what cannot be eliminated can still be reduced.

He then turns to the claim that science rests on faith too. Expecting the Moon to rise again and expecting a dragon to cure cancer are both uncertain, he says, but not equally. He ends by suggesting that people who answer every improvement with ``but it's still gray'' have made peace with their own mistakes.

\responsehead

The essay has one idea: uncertainty comes in degrees. The idea is correct and old. Asimov made the same argument in 1989, and the essay credits him in its last line, through a commenter. The nirvana fallacy and false equivalence already named the error. The essay's contribution is a new name and a set of opponents.

The opponents are the problem. The Sophisticate is a character in a novel. The man who shrugged off long odds was answering a figure the author admitted, the day before, he could not source. The bloggers are unnamed. The one real opponent, Tyler Cowen, is misquoted and misread. He argued that economic theory carries its own biases and that judgment is needed to correct them, which is a claim about which gray is lighter, the essay's own subject. The essay treats him as someone who thinks all grays are the same. The ``faithist'' is invented, given an angry tone, and mocked for it.

The examples are rigged in the same way. Gandhi against Stalin, green cheese against hydrogen, the Moon against an invisible dragon: every comparison is between extremes that no reader would confuse. The essay's thesis matters only where shades are close, and there it gives no help at all. It never says how to tell which of two nearby grays is lighter, which is the only hard question its thesis raises. It takes the side of the author's co-blogger, whose ``straightforward'' method had just produced an endorsement of a tax on height, without asking how light that gray is.

The essay also supplies a ready reply to critics. It was written for a blog named \textsc{Overcoming Bias}, and its first target is the objection that no one can overcome bias. By the end, anyone who says ``you're gray too'' without an itemized list is diagnosed with ``a devil's bargain with their own mistakes.'' The essay itself names no dark spots in its opponents' views. It names motives. A reader comes away with a label to attach to critics and a psychological story about why they criticize, which is the use of bias-talk the author had warned against nine months earlier.

On the one question with real content, whether science's reliance on induction is a kind of faith, the essay answers with the Moon's track record, which assumes induction. The author conceded the circularity six months later. The version in the collection is the dismissal.

In short: an old idea, credited last, defended against opponents who are fictional, unnamed or misquoted, and turned at the end into a diagnosis of the people who disagree.
\end{afterword}
